    \SetPalette{fancy}
    \sectionTitle{Coluna 3: Paleta Fancy}
    
    \begin{mainbox}[center title]{Check out our GitHub Repository \faGithub}
    \begin{center}
        \begin{minipage}[c]{0.2\linewidth}
            \centering
            \qrcode[height=\linewidth]{https://github.com/timotheosf}
        \end{minipage}
    \end{center}
    \end{mainbox}

    \PrintPalette
    \vspace{1cm}

    Aqui podemos testar uma lista para ver a cor mudando:
    \begin{itemize}
        \item Primeiro item de teste.
        \item Segundo item com \Highlight{destaque no texto}.
    \end{itemize}

    \begin{mainbox}{A caixa \texttt{mainbox}}
        A função de partição está ligada à energia livre:
        \begin{equation}
            F(T,V,N)=-k_BT\ln\mathcal{Z}_C.
        \end{equation}
    \end{mainbox}

    \begin{highlightbox}{A caixa \texttt{highlightbox}}
        Testando o visual secundário.
    \end{highlightbox}

    \begin{darkbox}{A caixa \texttt{darkbox}}
        Ideal para fundos densos.
    \end{darkbox}

    \begin{graybox}{A caixa \texttt{graybox}}
        Para comentários paralelos.
    \end{graybox}